% GENERATED FILE. Edit canonical lessons, then run npm run generate:books.
% canonical-source-hash: fnv1a64:fdaf1c716fc22606
% canonical-lessons: ML-C416-niti, ML-C416-aniti, ML-C416-anyayamaya, ML-C416-manusyavakasam, ML-C416-vivecanam

\chapter{Justice and Rights}
\label{ch:ml-justice-and-rights}
\hlchaptermodality{416}

\begin{chapteropening}
\textbf{By the end of this chapter:} \emph{I can say justice, injustice, unfair, human rights and discrimination.}

Complete the last lesson of chapter 416: I can say justice, injustice, unfair, human rights and discrimination.
\end{chapteropening}

\section[\texorpdfstring{\ml{നീതി} (nīti)}{നീതി (nīti)}]{\ml{നീതി} (nīti) — justice}

\label{lesson:ML-C416-niti}



\begin{quote}
Before the new one: say the Malayalam for equality, then the Malayalam for freedom.
\end{quote}

\subsection*{You'll want to know: \ml{നീതി}}
\textbf{\ml{നീതി}} — \emph{nīti} — \textquotedblleft{}justice\textquotedblright{}.

\begin{quote}
\ml{ആ} \ml{കുടുംബത്തിന്} \ml{ഇനിയും} \ml{നീതി} \ml{കിട്ടിയില്ല}.
\end{quote}

\begin{quote}
\emph{ā kuṭumbattinŭ iniyuṁ nīti kiṭṭiyilla.}
\end{quote}

\begin{quote}
That family has still not got justice.
\end{quote}

Justice is \ml{കിട്ടുക}, received, or \ml{നിഷേധിക്കുക}, denied.

\begin{cousinweb}
\textbf{\ml{നീതി}} (\emph{nīti}) is borrowed from Sanskrit \emph{nīti}, \textquotedblleft{}leading, guidance; conduct\textquotedblright{} (CDIAL 7546; Monier-Williams), from the root \emph{nī}, \textquotedblleft{}to lead\textquotedblright{}.
\end{cousinweb}

\begin{grammarlens}[title={What you've built}]
One more word for giving reasons and opinions.
\end{grammarlens}

\subsection*{Your turn}
\begin{itemize}
\raggedright
  \item \emph{Say it:} \emph{nīti}
  \item \emph{Say it:} \emph{nīti}, once more
  \item \emph{Say it:} \emph{ā kuṭumbattinŭ iniyuṁ nīti kiṭṭiyilla.} — That family has still not got justice.
  \item \emph{Recall:} say the Malayalam for equality, then the Malayalam for freedom, then say \emph{nīti} again
\end{itemize}

\begin{tcolorbox}[breakable,colback=teal!4,colframe=teal!35!black,title={Before you move on}]
What does \ml{നീതി} mean? (\textquotedblleft{}Justice\textquotedblright{}.) Say it once more.
\end{tcolorbox}
\section[\texorpdfstring{\ml{അനീതി} (anīti)}{അനീതി (anīti)}]{\ml{അനീതി} (anīti) — injustice}

\label{lesson:ML-C416-aniti}



\begin{quote}
Before the new one: say the Malayalam for freedom, then the Malayalam for justice.
\end{quote}

\subsection*{You'll want to know: \ml{അനീതി}}
\textbf{\ml{അനീതി}} — \emph{anīti} — \textquotedblleft{}injustice\textquotedblright{}.

\begin{quote}
\ml{അനീതി} \ml{കണ്ടാൽ} \ml{അവൾ} \ml{മിണ്ടാതിരിക്കില്ല}.
\end{quote}

\begin{quote}
\emph{anīti kaṇṭāl avaḷ miṇṭātirikkilla.}
\end{quote}

\begin{quote}
If she sees injustice, she will not stay silent.
\end{quote}

The Sanskrit prefix \ml{അ}- or \ml{അൻ}- makes an opposite: \ml{നീതി}, \ml{അനീതി}.

\begin{cousinweb}
\textbf{\ml{അനീതി}} (\emph{anīti}) is borrowed from Sanskrit \emph{anīti}, \textquotedblleft{}impropriety, injustice\textquotedblright{} (Monier-Williams): \emph{a-}, \textquotedblleft{}not\textquotedblright{}, and \emph{nīti}.
\end{cousinweb}

\begin{grammarlens}[title={What you've built}]
One more word for giving reasons and opinions.
\end{grammarlens}

\subsection*{Your turn}
\begin{itemize}
\raggedright
  \item \emph{Say it:} \emph{anīti}
  \item \emph{Say it:} \emph{anīti}, once more
  \item \emph{Say it:} \emph{anīti kaṇṭāl avaḷ miṇṭātirikkilla.} — If she sees injustice, she will not stay silent.
  \item \emph{Recall:} say the Malayalam for freedom, then the Malayalam for justice, then say \emph{anīti} again
\end{itemize}

\begin{tcolorbox}[breakable,colback=teal!4,colframe=teal!35!black,title={Before you move on}]
What does \ml{അനീതി} mean? (\textquotedblleft{}Injustice\textquotedblright{}.) Say it once more.
\end{tcolorbox}
\section[\texorpdfstring{\ml{അന്യായമായ} (anyāyamāya)}{അന്യായമായ (anyāyamāya)}]{\ml{അന്യായമായ} (anyāyamāya) — unfair; unjust}

\label{lesson:ML-C416-anyayamaya}



\begin{quote}
Before the new one: say the Malayalam for justice, then the Malayalam for injustice.
\end{quote}

\subsection*{You'll want to know: \ml{അന്യായമായ}}
\textbf{\ml{അന്യായമായ}} — \emph{anyāyamāya} — \textquotedblleft{}unfair; unjust\textquotedblright{}.

\begin{quote}
\ml{അത്} \ml{തികച്ചും} \ml{അന്യായമായ} \ml{ഒരു} \ml{തീരുമാനമായിരുന്നു}.
\end{quote}

\begin{quote}
\emph{atŭ tikaccuṁ anyāyamāya oru tīrumānamāyirunnu.}
\end{quote}

\begin{quote}
That was a completely unfair decision.
\end{quote}

The noun is \ml{അന്യായം}, a wrong.

\begin{cousinweb}
\textbf{\ml{അന്യായമായ}} (\emph{anyāyamāya}) is built from Sanskrit \emph{a-}, \textquotedblleft{}not\textquotedblright{}, and \emph{nyāya}, \textquotedblleft{}rule, right method\textquotedblright{} (CDIAL 7614), with the adjective ending -\ml{ആയ}.
\end{cousinweb}

\begin{grammarlens}[title={What you've built}]
One more word for giving reasons and opinions.
\end{grammarlens}

\subsection*{Your turn}
\begin{itemize}
\raggedright
  \item \emph{Say it:} \emph{anyāyamāya}
  \item \emph{Say it:} \emph{anyāyamāya}, once more
  \item \emph{Say it:} \emph{atŭ tikaccuṁ anyāyamāya oru tīrumānamāyirunnu.} — That was a completely unfair decision.
  \item \emph{Recall:} say the Malayalam for justice, then the Malayalam for injustice, then say \emph{anyāyamāya} again
\end{itemize}

\begin{tcolorbox}[breakable,colback=teal!4,colframe=teal!35!black,title={Before you move on}]
What does \ml{അന്യായമായ} mean? (\textquotedblleft{}Unfair; unjust\textquotedblright{}.) Say it once more.
\end{tcolorbox}
\section[\texorpdfstring{\ml{മനുഷ്യാവകാശം} (manuṣyāvakāśaṁ)}{മനുഷ്യാവകാശം (manuṣyāvakāśaṁ)}]{\ml{മനുഷ്യാവകാശം} (manuṣyāvakāśaṁ) — human rights}

\label{lesson:ML-C416-manusyavakasam}



\begin{quote}
Before the new one: say the Malayalam for injustice, then the Malayalam for unfair.
\end{quote}

\subsection*{You'll want to know: \ml{മനുഷ്യാവകാശം}}
\textbf{\ml{മനുഷ്യാവകാശം}} — \emph{manuṣyāvakāśaṁ} — \textquotedblleft{}human rights\textquotedblright{}.

\begin{quote}
\ml{മനുഷ്യാവകാശം} \ml{എല്ലാവരുടേതുമാണ്}.
\end{quote}

\begin{quote}
\emph{manuṣyāvakāśaṁ ellāvaruṭētumāṇŭ.}
\end{quote}

\begin{quote}
Human rights belong to everyone.
\end{quote}

Also in the plural: \ml{മനുഷ്യാവകാശങ്ങൾ}.

\begin{cousinweb}
\textbf{\ml{മനുഷ്യാവകാശം}} (\emph{manuṣyāvakāśaṁ}) is a Malayalam compound built from Sanskrit \emph{manuṣya}, \textquotedblleft{}human\textquotedblright{} (CDIAL 9828), and \ml{അവകാശം}, \textquotedblleft{}a right\textquotedblright{}, from Sanskrit \emph{avakāśa}, \textquotedblleft{}room, opportunity\textquotedblright{} (CDIAL 728).
\end{cousinweb}

\begin{grammarlens}[title={What you've built}]
One more word for giving reasons and opinions.
\end{grammarlens}

\subsection*{Your turn}
\begin{itemize}
\raggedright
  \item \emph{Say it:} \emph{manuṣyāvakāśaṁ}
  \item \emph{Say it:} \emph{manuṣyāvakāśaṁ}, once more
  \item \emph{Say it:} \emph{manuṣyāvakāśaṁ ellāvaruṭētumāṇŭ.} — Human rights belong to everyone.
  \item \emph{Recall:} say the Malayalam for injustice, then the Malayalam for unfair, then say \emph{manuṣyāvakāśaṁ} again
\end{itemize}

\begin{tcolorbox}[breakable,colback=teal!4,colframe=teal!35!black,title={Before you move on}]
What does \ml{മനുഷ്യാവകാശം} mean? (\textquotedblleft{}Human rights\textquotedblright{}.) Say it once more.
\end{tcolorbox}
\section[\texorpdfstring{\ml{വിവേചനം} (vivēcanaṁ)}{വിവേചനം (vivēcanaṁ)}]{\ml{വിവേചനം} (vivēcanaṁ) — discrimination}

\label{lesson:ML-C416-vivecanam}



\begin{quote}
Before the new one: say the Malayalam for unfair, then the Malayalam for human rights.
\end{quote}

\subsection*{You'll want to know: \ml{വിവേചനം}}
\textbf{\ml{വിവേചനം}} — \emph{vivēcanaṁ} — \textquotedblleft{}discrimination\textquotedblright{}.

\begin{quote}
\ml{ജോലിയിൽ} \ml{സ്ത്രീകളോട്} \ml{വിവേചനം} \ml{പാടില്ല}.
\end{quote}

\begin{quote}
\emph{jōliyil strīkaḷōṭŭ vivēcanaṁ pāṭilla.}
\end{quote}

\begin{quote}
There must be no discrimination against women at work.
\end{quote}

Against a group: -\ml{ഓട്} \ml{വിവേചനം}.

\begin{cousinweb}
\textbf{\ml{വിവേചനം}} (\emph{vivēcanaṁ}) is borrowed from Sanskrit \emph{vivecana}, \textquotedblleft{}discriminating, distinguishing\textquotedblright{} (Monier-Williams). The Sanskrit word is neutral, telling things apart; Malayalam uses it for treating people unequally.
\end{cousinweb}

\begin{grammarlens}[title={What you've built}]
One more word for giving reasons and opinions.
\end{grammarlens}

\subsection*{Your turn}
\begin{itemize}
\raggedright
  \item \emph{Say it:} \emph{vivēcanaṁ}
  \item \emph{Say it:} \emph{vivēcanaṁ}, once more
  \item \emph{Say it:} \emph{jōliyil strīkaḷōṭŭ vivēcanaṁ pāṭilla.} — There must be no discrimination against women at work.
  \item \emph{Recall:} say the Malayalam for unfair, then the Malayalam for human rights, then say \emph{vivēcanaṁ} again
\end{itemize}

\begin{tcolorbox}[breakable,colback=teal!4,colframe=teal!35!black,title={Before you move on}]
What does \ml{വിവേചനം} mean? (\textquotedblleft{}Discrimination\textquotedblright{}.) Say it once more.
\end{tcolorbox}
